\section{Problem Formulation}
\label{sec:problem}

\subsection{Notation}

Let $\fileset$ denote the set of source files an agent may interact
with, drawn from a fixed set of languages $\langset = \{\text{Go},
\text{Python}, \text{JavaScript}, \text{TypeScript}, \dots\}$. For a
file $\file \in \fileset$ we write $\loc(\file)$ for its line count
and $\tokensize(\file)$ for its token count under the LLM's tokenizer.
We assume $\tokensize(\file)$ is bounded by $c_\text{tok} \cdot \loc(\file)$
for a per-language constant $c_\text{tok} \in [8, 14]$ on realistic
source code; on the Blackrim test corpus the empirical mean is
$\approx 10$ tokens per line (\cref{sec:eval}).

Let $\outline : \fileset \to \fileset$ denote the canonical outline
operator: given a file $\file$, $\outline(\file)$ is a Markdown
summary of its top-level declarations, package documentation, imports,
and per-language significant nested constructs. We will write
$|\outline(\file)|$ for the outline's token count and define the
\emph{compression ratio}
\[
  \ratio(\file) \;\triangleq\; \frac{|\outline(\file)|}{\tokensize(\file)}.
\]
The smaller $\ratio$, the more compressive the outline.

Let $\budget \in \mathbb{N}$ denote a target token budget per outline
(default $\budget = 300$). The outline operator MUST yield
$|\outline(\file)| \leq \budget$ via the truncation precedence in
\cref{sec:algorithm:truncation}; conformance is an empirical claim
measured in \cref{sec:eval}.

Let $\threshold \in \mathbb{N}$ denote the LoC threshold above which a
\texttt{Read} should be preceded by an outline call (default
$\threshold = 200$).

An \emph{agent} $\agent$ takes a \emph{task} $\task$ and produces a
\emph{success} indicator $\success(\task; \agent) \in \{0, 1\}$
observable through one of three channels (\cref{sec:problem:success}).

\subsection{The read-tax reduction problem}

For a population of agent--file--task triples $(\agent, \file,
\task)$, the read-tax minimisation problem is:
\[
  \min_{\pi_\text{read}}\;
    \mathbb{E}_{(\agent, \file, \task)}\!
      \big[\, \cost_\text{read}(\agent, \file, \task; \pi_\text{read}) \,\big]
  \quad
  \text{s.t.}
  \quad
  \mathbb{E}\!\big[\, \success(\task; \agent, \pi_\text{read}) \,\big]
    \;\geq\; \success_0 - \epsilon,
\]
where $\cost_\text{read}$ is the input-token cost of the agent's
file-reading activity under reading policy $\pi_\text{read}$,
$\success_0$ is the baseline success rate when the agent reads
files in full, and $\epsilon$ is a quality-loss tolerance. The
outline-first reading policy $\pi^\star_\text{outline}$ is the
proposal:
\[
  \pi^\star_\text{outline}(\file) \;=\;
    \begin{cases}
      \text{full \texttt{Read}} & \text{if } \loc(\file) < \threshold, \\
      \outline(\file) \text{ first, then full \texttt{Read} \emph{iff} needed}
        & \text{if } \loc(\file) \geq \threshold.
    \end{cases}
\]
The agent retains the right to fetch the body --- the policy
encourages the outline as the first lens, not the only lens.

\subsection{The false-negative-only contract for symbolic surgery}
\label{sec:problem:falsenegative}

The write-side problem has a different shape. Let
$\intentset$ denote the registered set of transform intents
(\texttt{rename-symbol}, \texttt{extract-to-package},
\texttt{fix-imports}, \texttt{remove-console}, etc.) and let an agent
issue an intent $\intent \in \intentset$ on file $\file$. The
plan/execute pair (\cref{sec:system:planexecute}) produces a candidate
diff $\diff(\intent, \file)$ which is then submitted to a compile gate
$\gate$.

\begin{definition}[Gate verdict]
\label{def:gate}
$\gate(\diff, \file) \in \{\text{accept}, \text{reject}\}$ is
defined as $\text{accept}$ if the language's native syntax checker
parses $\file + \diff$ without error, and $\text{reject}$ otherwise.
\end{definition}

\begin{definition}[False-negative-only contract]
\label{def:false_negative}
A symbolic-surgery system satisfies the false-negative-only contract
if for every diff $\diff$:
\[
  \gate(\diff, \file) = \text{accept} \;\implies\;
    \diff \text{ does not corrupt valid programs.}
\]
The contract permits the system to reject diffs that would in fact be
safe (false negatives are allowed), but forbids accepting diffs that
silently break the program (false positives are forbidden).
\end{definition}

The contract is intentionally asymmetric. A false negative --- the
gate rejects a refactor that would have worked --- is an inconvenience
the agent can recover from (it tries again, perhaps with a smaller
scope). A false positive --- the gate accepts a refactor that breaks
the program --- is the failure mode we cannot tolerate, because it
poisons the downstream task with a quietly-broken codebase.

\subsection{Three observable signals for $\success$}
\label{sec:problem:success}

In production we observe the success indicator $\success$ through
three channels of varying fidelity:

\begin{enumerate}[leftmargin=*]
  \item \textbf{Compile / lint / test signal.} The most direct: did
        the post-edit codebase pass its own build? This is available
        for compiled-language refactors via the LSP's compile model
        and for any project with a CI configuration.
  \item \textbf{Reviewer signal.} A human reviewer (or in
        \blackrim, a Reviewer worker agent
        \cite{blackrim_paper_repo}) marks the agent's output as
        accepted or rejected. This is the source of truth for tasks
        whose ``success'' is not mechanically checkable (e.g., a
        documentation edit, an ADR review).
  \item \textbf{Downstream-task closure.} The bd~(beads) issue the
        agent was working on closes successfully, with no follow-up
        \texttt{redo} issue filed within $T_\text{redo}$ (default 7
        days). This is noisy but cumulative: drift in the closure
        rate is detectable even when individual signals are missing.
\end{enumerate}

The three channels are not independent --- a refactor that fails
compile will also be rejected by the reviewer and will produce a
follow-up issue --- but they fail independently. \Cref{sec:eval} uses
all three.

\subsection{Assumptions}

We make four assumptions, each minimal and each empirically checked
later:

\begin{assumption}[Stable per-language significance]
\label{assn:significance}
For a fixed tree-sitter grammar version, the set of significant nested
constructs per language (\cref{sec:algorithm:significance}) is stable
across files: the same node-type names denote the same syntactic
constructs.
\end{assumption}

\Cref{sec:discussion:treesitter} discusses what happens when
\cref{assn:significance} fails because of a grammar version bump.

\begin{assumption}[Token-cost monotonicity in LoC]
\label{assn:tokens}
$\tokensize(\file)$ is approximately linear in $\loc(\file)$ with a
per-language constant $c_\text{tok} \in [8, 14]$. Files violating this
(e.g., enormous one-line generated arrays) are out of scope; the
\texttt{outline:skip} comment escape hatch
(\cref{sec:system:discipline}) is the mitigation.
\end{assumption}

\begin{assumption}[Native-syntax-check fidelity]
\label{assn:gate}
The compile gate $\gate$ uses the language's native parser
(\texttt{go vet}, \texttt{python -m py\_compile}, \texttt{tsc
--noEmit}, etc.), which we take to be ground truth for syntactic
correctness. Semantic correctness (type-checking) is a stronger gate
available only for typed languages; the Tier-1 LSP refactors enforce
it where available.
\end{assumption}

\begin{assumption}[Bounded adversarial input]
\label{assn:adversarial}
Outlines are emitted over source files an agent has been directed to
read. Vendored third-party code may contain hostile content (e.g., a
package doc with prompt-injection text); we mitigate this with
sanitisation (\cref{sec:algorithm:sanitization}) but do not promise
robustness against arbitrary adversarial inputs.
\end{assumption}
